\documentclass[12pt,a4paper,onecolumn]{article}
\title{Medical Trivia of Grey's Anatomy}
\author{}

% \definecolor{dark-green}{rgb}{0,0.25,0}
% \usepackage[utf8]{inputenc}

\usepackage{fontspec,xunicode,xltxtra}
\setmainfont[Ligatures=Required,Mapping=tex-text]{Georgia}
% \setmainfont[Ligatures=Required,Mapping=tex-text]{Palatino Linotype}
\setsansfont[Ligatures=Required,Mapping=tex-text]{Myriad Pro}
\setmonofont[Ligatures=Required,Mapping=tex-text]{Courier New}
\usepackage{xeCJK}
\setCJKmainfont{宋体}
\setCJKmonofont{黑体}

% \usepackage{indentfirst}


\usepackage[dvips,dvipsnames,svgnames]{xcolor}
\definecolor{light-gray}{gray}{0.95}

%生成PDF的链接
\usepackage[dvipdfmx]{hyperref}
\hypersetup{
    bookmarks=true,         % show bookmarks bar?
    bookmarksopen=true,
    pdfpagemode=UseNone,    % options: UseNode, UseThumbs, UseOutlines, FullScreen
    pdfstartview=FitH,
    pdfborder=1,
    pdfhighlight=/P,
    pdfauthor={wuxch},
    unicode=true,           % non-Latin characters in Acrobat’s bookmarks
    colorlinks,             % false: boxed links; true: colored links
    linkcolor=blue,         % color of internal links
    citecolor=blue,        % color of links to bibliography
    filecolor=magenta,      % color of file links
    urlcolor=cyan           % color of external links
}
% 使用目录
\usepackage{makeidx}
% 使用图形
\usepackage{graphicx}


\usepackage{titletoc}
\dottedcontents{section}[1.5em]{}{2.3em}{1pc}

% % 页眉页脚的设定
% \usepackage{fancyhdr}
% \pagestyle{fancy}
% \lhead{\includegraphics[height=0.2in]{zte_logo}}

% \usepackage{footmisc}
\usepackage{pfnote}
\usepackage{fnpos}

% % 嵌入的代码显示
% \usepackage{listings}
% \lstset{language=C++, breaklines, extendedchars=false}
% \lstset{keywordstyle=\color{blue},
%         commentstyle=\color{SeaGreen},
%         stringstyle=\ttfamily,
%         showstringspaces=false,
%         backgroundcolor=\color{light-gray}}
% \lstset{basicstyle=\small}

\usepackage{booktabs}

\usepackage{tipa}
\usepackage{phonetic}
\usepackage[gen]{eurosym}

% \usepackage{CJK}
\usepackage{makeidx}
% \usepackage{glossary}

\usepackage{fancyhdr}

\usepackage{lastpage}
\cfoot{\textsf{\thepage\ of \pageref*{LastPage}}}

\pagestyle{fancy}
\fancyhead{}
\renewcommand{\headrulewidth}{0pt}
\pagenumbering{arabic}

\usepackage{graphicx}
\usepackage{picinpar}

\setlength{\parindent}{0in}
\setlength{\parskip}{2.5ex plus4mm minus3mm}
% \setlength{\parskip}{.3\baselineskip}

\setcounter{section}{0}

% bibtex
\usepackage[super]{natbib}
% \usepackage[globalcitecopy,labelstoglobalaux]{bibunits}
% \makeatletter
% \renewcommand{\@bibunitname}{\jobname.\the\@bibunitauxcnt}
% \makeatother

% \renewcommand{\refname}{缩略语}
% \renewcommand{\refname}{}

\makeatletter                   % Reference list option change
\renewcommand\@biblabel[1]{#1:  } % from [1] to 1: 
\makeatother                    %

\newcommand{\citenumfont}[1]{\textsf{#1}}

% \usepackage{tipa}
% \usepackage{phonetic}

\newcommand{\phonetic}[1]{{\fontspec[Colour=CC00CC]{Microsoft Sans Serif} #1}}


%%% Local Variables:
%%% TeX-master: "medical_trivia"
%%% End:

\makeindex
\begin{document}
\date{}
\maketitle
\tableofcontents
\pagebreak
% \begin{sffamily}
\section*{Season 1}
\section*{Season 2}
\section*{Season 3}
\section{Season 4}
\subsection{A Deer in the Headlights: Holly the Deer Comes to Seattle Grace}

In the season premiere, A Change is Gonna Come, Izzie Stevens treats a patient whose anatomy she
never studied, a black-tailed deer played by Hollywood deer sensation Holly Berry.

Holly was hit by a car. That initial collision sets off a chain reaction of victims, including a
pregnant woman whose arm is severed in the accident, our
doe\cite{doe} is presumed dead and loaded into the back of a pick-up truck.
When the driver finds the severed arm and brings it to Seattle Grace, it's his son, Joey whose
heart-felt pleas convince Izzie to save the life before her, the life of a deer.

Last year approximately 3000 deer were killed in the State of Washington by automobiles. But
sometimes, they survive the collision and are cared for by one of the large animal
veterinarians\cite{veterinary} we spoke to when researching this story.

When a deer is hit by a car, and survives, the first thing that happens is that the deer goes into
shock, and to an untrained person, could appear to be dead as was the case in our episode.

When a deer goes into shock, it lays perfectly still with what our vets described as an incredibly
intense, wide-eyed look (yes, the ``deer in the headlights'' description is accurate) and has short
shallow breaths. Though it doesn't usually make a sound, the deer is in immediate need of medical
assistance.

Deer are ruminants\cite{ruminant}, which mean that they have a digestive system that is similar to a
cow or a horse. Their stomach is divided into four sections and they regurgitate\cite{regurgitate}
their cud\cite{cud}. Deer have to be careful, because if they lay on their side for too long they
can choke to death. This happens when the green partially digested food travels up the throat.
Another life threatening complication in the treatment of deer is the main stomach swelling. It
fills with gas, compromising the deer's ability to breathe. Vets expel the air by stabbing two
really, really big needles (16-gauge needles) into the abdomen, which lets gas out of the stomach.

If a deer is in severe shock, or if the pressure of the collapsed lung builds to a critical point,
the deer could go into v-fib, or a fatal heart rhythm. By using the same
defibrillators\cite{defibrillator} used on humans, a vet can shock the heart back into a regular
rhythm.

Deer have a highly developed adrenal\cite{adrenal} response to stress. This is called the ``fight or
flight'' response. When they come out of shock, they kick off the exam tables and try to get out of
whatever enclosed space they are in. Whenever possible, vets administer Valium or another
sedative\cite{sedative} to try to diminish this reaction, which is potentially dangerous for the
deer and the vet.

Though our deer at Seattle Grace came out of shock peacefully, don't try this at home! (Or at the
side of the road.) Deer are very strong and have incredibly sharp hooves. And there are cases of
well- meaning animal lovers who loaded injured deer into their car only to have a full sized buck
kicking out the windshield en route to a veterinary hospital.

All of the vets we spoke to had heartwarming stories of deer they had rescued. Sometimes a female
deer may leave a fawn\cite{fawn} behind. There are a number of rehabilitation centers that can raise
a fawn to maturity and work with it to reintroduce it to the wild.

Holly Berry, our actress-deer was rescued as a fawn herself. She went from wandering in the woods
without a mother to her glamorous life of being the most sought after deer in tinsel town.

\subsection{No Title}

Maybe you've looked in the mirror and thought; I could use a haircut. For someone with Body
Integrity Identity Disorder, their thoughts run all the way to, I really need to loose one of my
arms.

In Episode 405, Haunt You Everyday, James Miller, suffering from BIID, wants his left foot
amputated\cite{amputate}. Bailey says no way! So, James steals a power saw and goes for it himself.
Ouch!

Why would anyone want to get rid of a perfectly healthy foot? Because they believe, to the depths of
their heart, that their true self only has one foot. Or no feet at all! BIID sufferers are obsessed
with the idea that there given body is wrong and the body they see in their head, the one missing a
foot is the right body. And they will go to great lengths to force an amputation. It's not just feet
- arms, legs, even hands. There have been reports of people laying their arms on a railroad tracks,
burning their legs with dry ice and crushing their own hands.

BIID folk claim that they are not suffering a mental illness at all. It's like waking up with a
third arm they just want something they don't think is right, removed. They claim that their
condition is just like gender identity disorder (when someone is born a girl and believes they are
really a boy or vise versa). If someone can change their gender with surgery, why can't they become
an amputee?

Well, to have elective amputation, you have to find a doctor who is willing to do it. Putting
someone through dangerous surgery to lose a healthy limb puts the most revered medical principle,
first, do no harm, to the test. Enraging the medical community worldwide, one doctor in Scotland
performed a double amputation on a BIID patient because he thought not amputating was doing harm!
His reasoning was that the pain and suffering of living in body you don't want to be in is worse
that going through major surgery to live the rest of your life in a disabled state. Yeah, RIGHT!

There have been some attempts at treatment. Under the care of a psychiatrist, anti-depressants and
cognitive\cite{cognitive} therapy can help take away the obsessive thinking part of this baffling
disorder. BIID shares characteristics with Obsessive Compulsive Disorder and Body Dismorphia. There
are even theories that this is a neurological problem in the cerebral\cite{cerebral}
cortex\cite{cortex}. Damage in this area could mean the brain fails to recognize the limb as part of
the body. Either way, the condition is a big challenge for doctors and patients.

So, the next time you look in the mirror and you get a flood of thoughts about losing a few pounds,
be thankful it's not a few limbs!

\section{Season 5}
\subsection{Dream a Little Dream of Me }
\begin{flushright} --- Hypothermia
\end{flushright}

Being cold really sucks, but sometimes, it's just the thing to save your life.

How can hypothermia\cite{hypothermia} save lives? Everyone's heard stories of people trapped on
mountain tops or snow storms or at the South Pole who tragically die from the cold conditions.
Hypothermia causes shivering, blue lips, frostbite and eventually death.

But, sometimes people didn't die from being too cold. So, scientists began to study hypothermia
because every once in a while people (like Meredith) who have been trapped under cold water and
deprived of oxygen for up to twenty minutes miraculously survive. Scientists wanted to find out what
causes these unbelievable survivals. Could they recreate this on purpose to save a patient? They
experimented with therapeutic\cite{therapeutic} temperature management and found that cooling a
patient's core temperature can help reduce organ damage and prevent paralysis\cite{paralysis} or
even death. Therapeutic hypothermia is becoming a treatment option for heart attacks, strokes and
traumatic spinal cord and brain injuries.

One of the ways therapeutic temperature management works is by reducing
inflammation\cite{inflammation}. Like when you ice a sprained ankle. But, there's one important
difference. The spinal\cite{spinal} cord is housed in rigid bone - the spine - which can't expand to
accommodate the swelling nerves. As the spinal cord become larger than the bone it is incased in
it's crushed and damaged even more. In fact, the inflammation can cause more trouble than the
initial injury, so preventing swelling significantly increases recovery. Hypothermia also slows down
metabolism\cite{metabolism}. For every degree the body temperature is lowered, metabolism slows 6\%.
This slows down cell death and halts the release of toxic-chemicals that follow a spinal cord
injury.

Doctors cool patients by rushing cold saline\cite{saline} into the blood supply through the femoral
artery in the thigh. They drape the patients' head, neck, chest and groin area with special core
cooling blankets. Body temperature is closely monitored, because even though hypothermia has many
potential beneficial effects, it is also dangerous for the patient. So, it's a delicate balance
between cooling the body enough to reap\cite{reap} the benefits, but not so much that the patient's
heart stops.

Callie makes the argument to cool the patient down before taking him in to surgery to decompress his
spine. When the patient wiggles\cite{wiggles} his toes at the end of the episode, they know that the
cooling was indeed a success. Though scientists don't yet understand all the nuances\cite{nuance} of
therapeutic hypothermia, doctors are becoming more and more open to cooling their spinal cord
trauma\cite{trauma} patients. Thankfully, in this case, Callie's experiment was worth the risk.


\subsection{What Is The Worst Pain You'll Ever Experience? }

In Episode 503, ``Here Comes The Flood'', Meredith feels the pain of being torn between her
boyfriend who wants her to kick out all her other roommates and her friends, who need a place to
live. Retaking the intern exam is George's pain. And nothing says pain to Cristina like NOT being on
the cardiothoracic surgery rotation. But, none of these compare to Derek's patient, Phillip
Patmore's pain. He has been living with an intense unrelenting headache for seven years!

When it was clear that a simple aspirin wasn't going to stop the Phillip's pain, doctors prescribed
narcotics\cite{narcotics}, anti-seizure medications, antidepressants, anti-psychotics. None of these
worked. Phillip's tried everything except having part of his brain surgically removed and that is
exactly what Dr.~McDreamy, pain specialist and neurosurgeon extraordinaire is about to do.

The procedure is called a bilateral cingulotomy. Derek opens into Phillip's skull, targets part of
his frontal lobe\cite{lobe} and removes two small pieces of his brain. It's so radical that it is
only done in cases of extreme and unmanageable pain or OCD. The risks are serious: seizures,
cognitive\cite{cognitive} defects and behavioral changes and no one knows exactly why this technique
works

But, right before the surgery, Lexie Grey, Meredith's younger sister with a photographic memory,
remembers an obscure condition of a pinched sinus\cite{sinus} nerve which is described as the worst
pain you'll ever know. Lexie makes a connection. What if Phillip's headache is actually caused by a
pinched sinus nerve?

Inside the nasal\cite{nasal} passages there are turbinates, which regulate the flow of air through
the nose. They filter out harmful particles and add moisture to the air and warm it before it
travels to the lungs. Sometimes a previous infection or a trauma causes a glitch and a nerve in
Phillip's case the ethmodial nerve gets trapped between one of these turbinates and the nasal
septum. It's called a ``contact point'' or ``trigger''. Because this condition is so rare, the most
difficult part of the diagnosis is thinking of it. Usually sinus headaches involve inflammation and
swelling of the sinuses. But in Phillips case he had no current sinus infection or swollen sinus
tissue, so doctors had not considered this as a possibility. Once it is suspected, there are two
easy ways to rule out or confirm the diagnosis. The first, while the pain is in evidence, put
anesthetic\cite{anesthetic} on it. If the headache goes away, then a trapped ethmodial nerve is
indeed the problem. The second way is to push on the trigger and if the patient screams in pain.
Either way a simple surgery to clear out the blockage that is putting pressure on the nerve will fix
the problem permanently.

While Lexie is sometimes the cause of Mereidth's pain, for Phillip Patmore, little Lexipedia saved
the day!

\subsection{How to Spice Up A Boring Day in the Clinic}

In Episode 504,``Brave New World,'' Cristina is confronted with her least favorite thing in the
world - a non-surgical, loud-mouthed, patient from Cristina's least favorite place in the world -
the clinic. Mrs. Borsokowski arrives complaining of a rash. Then she turns feverish,
itchy\cite{itchy}, uncomfortable and unhappy to see Dr.~Yang - whom she calls,``the mean one''.
Cristina, desperate to get into the OR, makes the one diagnosis that could land her a surgery:
necrotizing fasciitis. This is the deadly infection that rapidly consumes flesh. The only treatment
is to surgically removed the infected flesh and pray that you remain ahead of the infection. This
means, limb amputation or death! All of which makes it the perfect diagnosis to spice up a boring
day in the clinic for Cristina. Unfortunately, it only takes Mark Sloan one peek at the rash to
know, this isn't flesh eating bacteria. He prescribes a topical cream.

But Mrs. Borsokowski doesn't have just a rash. She has something far worse-- and just as dangerous
as necrotizing fasciitis - toxic\cite{toxic} epidermal\cite{epidermal} necrolysis\cite{necrolysis}
(or TEN syndrome). Cristina goes to dermatology\cite{dermatology} to get rash cream, while Lexie
watches as Mrs. Borsokowski's condition goes from annoying to life threatening!

TEN syndrome is most typically an extreme allergic reaction to taking sulfa\cite{sulfa} drugs,
anticonvulsants\cite{convulsant}, NSAIDs, antiretroviral\cite{retrovirus} medications or
corticosteroids. Those drugs can set off a series of reactions, fever, coughing,
malaise\cite{malaise}, rash\cite{rash} on the trunk of the body. Painful lesions\cite{lesions} erupt
on the mucosal membrane\cite{membrane}, like the inside of your mouth. These are very painful and
often cause their victim to stop eating. If the reaction isn't halted, as in Mrs. Borsokowski's
case, the patient will go into respiratory\cite{respiratory} distress as the
epidermis\cite{epidermis} (the top layer of skin) separates from the derma\cite{derma} (the
underlying layer of skin). This is referred to as Nikolsky's sign -- the skin can literally peel off
with the slightest touch. It's like a worst sunburn you have ever had times 100! The patient suffers
from skin loss like burn victims do, even though they have never been burned. Doctors believe this
is from a biochemical progression set off by the drug which causes cells to change and then die.
Mrs. Borsokowski lost 60\% of all of the skin on her body and will go to the burn unit where they
will debride\cite{debride} her (remove the dead and sloughing skin) and wash it to keep infection at
bay.

Looks like Cristina's day at the clinic wasn't boring at all! Good thing she was looking for
something deadly serious because that day-- she found it!

\subsection{There's No ``I'' in Team}

\begin{flushright} --- Domino Kidney Transplant
\end{flushright}

One of the questions I'm often asked is,``How has your job as Medical Researcher of \emph{Grey's
Anatomy} changed your attitudes toward medicine.'' The answer is, when I renewed my driver's license
this time, I signed up to be an organ donor. One thing \emph{Grey's Anatomy} (and Denny Duquette)
has taught me is there just aren't enough organs from deceased donors to go around. In 1954, doctors
began performing living organ transplants when a living person donates a kidney or part of their
liver to a loved one who is in need of an organ. While living organ donation hasn't solved the organ
shortage, it has definitely helped close the gap. It's also created an additional problem, what if
someone in your life is generous enough to give you a kidney but their kidney doesn't match? In
episode 505, There's No``I'' in Team, Dr.~Miranda Bailey creates an innovative program that provides
an answer to that\ldots The Domino Kidney Transplant.

You know how when you stand up dominos in a pattern and knock one over all of the dominoes fall over
in quick succession? That's just what Bailey did with the domino kidney transplants. She found
someone who needs a kidney and had someone in his or her life who was not a match, but was willing
to give up a kidney to help. Bailey found five other donor/recipient pairs in this same situation
and she paired them up one by one, until there were five recipients and six donors who matched and
are ready to go. What happened to the sixth kidney?

The United Network for Organ Sharing (UNOS) has developed a system of organ allocation that is fair.
Well, as fair as it can be when dealing with something as fundamentally unfair as life-threatening
illness. In cases of directed donation (like if your sister gives you a kidney) the
ethic\cite{ethic}s is straightforward. The donor has the right to decide who she gifts her organ to.
But, as soon as a licensed transplant center gets involved with the pairing process, the ethics
becomes more difficult. The rules say the person at the top of the list gets the next kidney.
Doctors don't allocate organs. UNOS does. Yet, Bailey knows six people who are willing to donate
organs in order for their loved ones to receive the organs. With such a grave\cite{grave} organ
shortage it seems a shame not to take them up on their generosity\cite{generosity}. So, this problem
is solved when one of the kidneys is given to the person at the top of the list. They become the
sixth recipient.

Who is the sixth donor? Where did they come from? Typically the sixth donor is what is called an
altruistic\cite{altruistic} donor. Someone who wants to donate an organ just to be generous. This
altruistic donor is actually the first domino that trips the chain reaction, because without the
unpaired kidney there would be no kidney to pair with the person at the top of UNOS's list.

Now for the question that everyone wants answered: Would the kidney that Meredith dropped really be
used? I'll answer that question with a question. Now that you know how rare a good kidney is and how
much work goes into pairing kidneys and recipients what do you think?

For more information on Kidney Transplants or becoming an organ donor go to: www.donatelife.net


\subsection{Life During Wartime}
\begin{flushright} --- Tumors
\end{flushright}

Bailey has a big problem. Richard told her to remove Tori Begler's abdominal tumor. Seems easy
enough. After all, that's one of the things surgeons are really good at, removing tumors. Only, it's
a little more complicated than that. Tori's tumor is deep down within her abdominal
cavity\cite{cavity}. So her stomach, intestines\cite{intestine}, spleen\cite{spleen}, kidneys and
liver are on top of the tumor and it's impossible to get at it. Meredith, who is looking at her
Anatomy Jane Doll, has an idea. What if you took out all of the organs one by one and removed the
tumor and then put all of the organs back in the body? The answer is that won't work. By the time
you took the last organ out of the body, the first one you took out would be dead from lack of
blood. However, Meredith's idea inspires the Chief. What if you left the organs attached to their
blood sources and took all of the abdominal organs out in one big chunk? Seems easy enough -- on the
plastic doll.

While it seems like an obvious solution it isn't. Typically, surgeons put their hands into the body
cavity to remove tumor, so the organs remain in place. You know how sometimes when you take one of
those 3D puzzles apart, it's hard to put back together again? Well, this is a living 3D puzzle,
which means, the pieces change size as you go. It also means that they are connected together by an
intricate\cite{intricate} system of arteries and veins, which supply life giving oxygen rich blood
to the organs. If one of these is disconnected incorrectly, Tori could bleed out. Conversely, if any
of the organs are disconnected from its source of blood for too long, it will die. To keep
everything in tact and in relation to everything else, they lift all of the organs out at the same
time, sort of like they are removing an engine from under the hood\cite{hood} of a car. Only in this
instance, the engine is still connected to the body by two blood vessels: the aorta\cite{aorta} and
the vena\cite{vena} cava and two digestive organs: the esophagus\cite{esophagus} and the large
bowel\cite{bowel}. This part goes well, now all they have to do is remove the tumor.

Tumors are little vampires\cite{vampire}. They create veins and arteries and siphon\cite{siphon} off
blood for themselves. The bigger the tumor, the more blood supply. This tumor is lemon sized, which
for a tumor, is enormous for a tumor in this location. Which means they have to figure out what to
do with all of the disconnected blood vessels once they remove the tumor. Since the tumor itself is
surrounded by blood vessels, once the tumor is removed all of the vessels will have to be repaired.
For this, Erica is using a kind of surgical mesh called gortex.

Meredith's inspiration and Richard's ingenuity\cite{ingenuity} combined with Erica Hahn's surgical
skills and Baily's preserverence save Tori. This ground breaking surgery will draw positive
attention to Seattle Grace and improve their standing in the eyes of their peers. Tori is in good
hands and is well on her way to recovery. In the end, both Anatomy Jane and Tori have all their
organs back where they belong and are expected to make a full recovery.

For more information about cancer please go to: \url{http://www.cancer.gov/}

\subsection{Rise Up}
\begin{flushright} --- Alcohol Septal Ablation\cite{ablation}
\end{flushright}

Michael Norris has been and Dr.~Hahn's patient a long, long time. He has a
congenital\cite{congenital} heart condition, hypertrophic\cite{hypertrophic}
cardiomyopathy\cite{cardiomyopathy} that makes him a candidate to receive a heart transplant. In
fact, he's been near the top of the heart transplant list since Season 2 of Grey's Anatomy. The
heart Denny received before he died would have been Mike's heart if the machine that was pumping
Denny's heart (the LVAD) had not mysteriously failed. Now, Mike Norris has yet to receive a match.
Time is running out. His heart is too weak to withstand open-heart surgery, if Dr.~Hahn does
nothing, Mike won't live long enough to get a new one. Lucky for Mike, Erica knows just the thing to
do, a cutting edge procedure called an alcohol septal ablation.

During alcohol ablation Erica will inject portions of Mike's heart with 100\% absolute alcohol, a
substance that is lethal\cite{lethal} to heart cells. Once the cells die, they shrink. As the heart
muscle expires, the thickened muscle becomes thinner scar tissue. Mike's blood will be able to flow
more freely, his heart will pump more efficiently and his lungs will be able to receive more oxygen.
In short, his quality of life will greatly improve.

There are four stages to this procedure. First, Erica threads a tiny temporary pacemaker into Mike's
heart. If the procedure goes well, they'll take it out the next day, but until then, it will ensure
that Mike's heart beats rhythmically. Secondly, Erica injects a contrast die into the vessels. This
illuminates them in x-rays and allows Erica to see exactly where she will inject the alcohol. Then,
Dr.~Hahn inserts a balloon-tipped catheter into a branch of the targeted artery. When inflated, it
blocks off the blood flow and limits the exposure of the alcohol so that it can safely be injected.
As, the alcohol is injected into the heart, Erica is causing a heart attach. Though Mike is sedated
and has sufficient amounts of pain medication, he feels discomfort during the procedure. Finally,
Erica performs another angiogram to check her work. It takes two tries and an appearance from Denny,
but Mike is better off for having come to Seattle Grace.

For more information about cardiomyopathy please go to:

\url{http://www.nhlbi.nih.gov}

\subsection{These Ties That Bind}
\begin{flushright} --- Autism\cite{autism} Spectrum Disorders
\end{flushright}

Bailey agrees with Alex that the new heart surgeon, Dr.~Virginia Dixon, is a bit off. She makes eye
contact a little too intensely, her speech pattern is strangely formal and she can't seem to bend
the rules, any rule, even a little bit. When Bailey and Dixon disagree about a patient's care,
Bailey undermines Dixon's authority as an Attending Physician rather than dealing with her strange
behavior. It turns out, Dr.~Dixon has a neurological condition called Asperger's syndrome. Dr.~Dixon
may seem like she doesn't have emotions, but it isn't the case. Dixon stands up for herself in the
end, and is a brave advocate for herself and her condition.

Asperger's is one of a group of developmental disorders known as autism spectrum disorders (ASD).
People with Asperger's have trouble fitting in to ``normal'' social interactions. They have
difficulty interpreting and sending social cues. They tend to take things literally. They often
stare a little too intensely, talk too loudly, have monotone intonations\cite{intonation} or use an
antiquated speech patterns. They adopt repetitive routines and rituals. They'll often become fixated
on one thing and study it intensely. For Dixon, her one passion is hearts.

No one knows what causes any of the autism spectrum disorders. They range in severity from disabled
to perhaps a little ``off'' by normal conventions, but truly gifted. Though an ASD may be diagnosed
as early as three years old, the eventual manifestation\cite{manifestation} of the disorder can
change dramatically throughout a person's life. Early intervention, changes in diet and nutrition,
some medications and theraputic socialization can all assist a child with ASD to connect to others
and gain life skills. Specialist intervention to help with at an early age significantly improves a
child with ASD's social and behavioral development, so if a child you know exhibits symptoms, it's
good to have a pediatrician do an exam.

During the course of working on this episode the research department stumbled upon the story of
Dr.~Temple Grandin. She's an autism expert (which is how Moira found her) and lectures all over the
globe. She's also one of the world's foremost animal behavior specialists. Her journey to understand
her own condition made her uniquely qualified to connect to animals and she has used her gift to
improve the conditions of livestock worldwide. She's a fascinating free thinker who makes me strive
to be a better smarter person. If you want a special treat, catch out some of her lectures online.
Oh, and here's a link in case you'd like more information about autism.

For more information on Autism Spectrum Disorders:

\url{http://www.cdc.gov/ncbddd/autism/overview.htm}

\subsection{In the Midnight Hour}
\begin{flushright}
--- Hypochondria
\end{flushright}
 
Alex's patient, MRS. HAMMER, has a pain in her stomach. The pain is so bad she thinks it must be
stomach cancer. Her husband thinks it's diarrhea\cite{diarrhea}. Turns out, she's a hypochondriac
who took antibiotics\cite{antibiotic} she ordered online to cure a self-diagnosed staph infection
that was actually a pimple\cite{pimple}. Unfortunately, the antibiotics killed all of the bacteria
in her body, including the good bacteria in her digestive tract that's necessary for digestion. So,
hypochondria ended up making her very sick.

To my knowledge there is no scientific study to determine the percentage of writers who are also
hypochondriacs. My guess is the percentage is somewhere around one hundred. At least, that's the
statistic in the Grey's Anatomy's writer's room. We're always asking Zoanne Clack, the resident
writer/physician, if some random obscure symptom we've imagined could possibly mean
immanent\cite{immanent} death. Like, last week I had a lump on my finger. I was sure it was flesh
eating bacteria and I was probably going to die if she wouldn't amputate\cite{amputate} my finger.
The symptoms where present, my skin was warm to the touch. Zoanne said my skin was warm to the touch
because I was just holding a cup of steaming coffee. She wrote me a note excusing me from doing
dermatology\cite{dermatology} research for the rest of the day, but refused to do the amputation.

I'm not alone. Millions of people suffer from hypochondria. According to some studies, around 3\% of
the US population is convinced they're sick. With the Internet, hypochondriacs are able to google
their symptoms and diagnose themselves. This is great for us, but doctors don't like it so much. The
largest complaint I hear from physicians is they spend a lot of time educating patients about drugs
on television commercials, telling their patients why they don't have some rare thing they saw on a
medical drama or explaining why they aren't ordering the new diagnostic test that was on the news
last night. Physicians have limited amounts of time with patients and prefer to educate them on
conditions, medications, diseases, etc. that actually have direct relevance to their patient's
health. I'm sympathetic to the plight\cite{plight} of the doctors but let's face it, cyberchondriacs
are here to stay and even before the easy access to medical information, hypochondriacs have been
self obsessing.

My favorite thing about hypochondria is the irony of our condition. We investigate every detail of
imagined disease and vague symptoms, but our true plight -- the hypochondria itself -- is a mere
afterthought or truth or dare confession for us. But, hypochonadrisis (the formal name) is a
psychiatric disorder and is defined in the DSM IV, the bible of psychology. It's considered a
somatoform disorder (having to do with complaints that result in a prescription) and is defined as,
``$[$a$]$ preoccupation with fears of having, or the idea that one has, a serious disease based on a
misinterpretation of one or more bodily signs or symptoms.'' It goes on to say, even when doctors
assure us we're fine, we're still convinced. It's often accompanied by other psychological disorders
such as clinical depression, OCD \& other phobias\cite{phobia} or anxiety disorders.

No one is sure what causes hypochondria. There are some schools of thought that link it to a
specific trauma or a traumatic illness in one's life. I can't think of anything in my life that
would qualify as such. But, my finger is hurting again, so I'm going to go and see if Zoanne has
changed her mind about performing that amputation.

For more information about hypochondria please go to:

% \url{http://www.nlm.nih.gov/medlineplus/ency/article/001236.htm}

\pagebreak

\bibliographystyle{d:/wuxch/emacs_home/vocabulary/wuxchvcb}
\bibliography{d:/wuxch/emacs_home/vocabulary/vocabulary}

% \end{sffamily}
\end{document}
